% ch11 —— 第 11 章 高橋篤史「特異点とミラー対称性」
% 本章原书节号是阿拉伯数字（1, 2, 3, …），与全书默认的罗马数字不同。
% 按 STYLE_GUIDE §3.2「各章按原书自己的节号体例」：章前局部改写，章末 reflist 之后恢复。
\renewcommand{\thesection}{\arabic{section}}
% ===== PDF p.145 (原书 p.137) =====
\chapter{奇点与镜像对称}
\chapterauthor{高橋篤史}{[大阪大学大学院理学研究科]}

\section{引言}

第 5 章讲述了从奇点导出拓扑场论的情形。本章想解说与同调镜像对称相关的话题。

所谓奇点，是指与周围的结构相比，性质格外不同、尤其是具有优异性质的点。更具体地说，某个图形的一部分图形坍缩成一点，该部分图形所拥有的信息被凝缩，那个特别的一点便带有许多对称性与整体结构的信息。

这里要注意的，是被称为超曲面奇点的、出现在由一个多项式的零点集合所表示的图形上的特殊点，也就是性质非常好的奇点。第 5 章也写过：本来光滑的非交换流形，正因为我们是用交换几何（复代数簇）去看它，它的性质才以奇异性的形式显现出来——可以这样理解。本章期待能传达这种氛围。

% ===== PDF p.146 (原书 p.138) =====
\section{Hodge 数与拓扑镜像对称}

那么，在观察奇点的镜像对称之前，先讲一讲 Calabi--Yau 流形的拓扑镜像对称。典范丛平凡的那种代数复流形称为 Calabi--Yau 流形。说得再通俗些，就是：第 5 章出现过的复射影空间 $\CP^n$ 的子空间 $X$，由有限个齐次多项式 $f_1(x_0,\dots,x_n),\dots,f_k(x_1,\dots,x_n)$ 的共同零点定义，并且第一 Chern 类 $c_1(X)$ 为零。就是这样的对象。典型的例子是 $\CP^n$ 中由 $n+1$ 次多项式的零点所表示的图形。
% [Tutor Note: 原书此处 $f_k$ 的变量印作 $x_1,\dots,x_n$，与前面的 $f_1(x_0,\dots,x_n)$ 不一致；按 $X\subset\CP^n$ 的齐次坐标应为 $x_0,\dots,x_n$。原书如此，照录。]

作为表示复流形图形特征的量，有称为 Hodge 数的整数
%
\[
  h^{p,q}(X):=\dim_{\C}H^{q}(X,\Omega_X^{p}).
\]
%
这是把“图形 $X$ 中含有多少部分图形”与“图形 $X$ 中有怎样的（高维的）洞”能简洁地表示出来的不变量。前面提到的复射影空间 $\CP^n$ 的情形，当 $p=q=i=0,\dots,n$ 时 $h^{i,i}(\CP^n)=1$，而其余的 $p,q=0,\dots,n$、$p\ne q$ 时全都 $h^{p,q}(\CP^n)=0$。Calabi--Yau 流形的特点是，按定义 $h^{\dim_{\C}X,0}(X)=1$ 成立。

Hodge 数带有若干对称性。例如，对复系数取复共轭这种对称性，会扩张为复流形 Hodge 数的对称性。它导出互换 $h^{p,q}(X)\leftrightarrow h^{q,p}(X)$。由于这一操作是图形的对称性，可以证明复流形的 Hodge 数满足等式
%
\[
  h^{p,q}(X)=h^{q,p}(X).
\]
%
此外，还有一种叫作 Serre 对偶的奇妙对称性，同样已知等式
%
\[
  h^{p,q}(X)=h^{dx-p,dx-q}(X),\qquad dx:=\dim_{\C}X
\]
%
成立。

看到这两种对称性，您不觉得“是不是也存在把 $h^{p,q}(X)$ 与 $h^{dx-p,q}(X)$ 互换的对称性呢？”吗？正是如此。引起这种剩下的 Hodge 数互换的对称性，正是镜像对称。不过，这次它不是同一个图形上的对称性，而是联系不同图形之间的对称性，这一点很有意思。也就是说，对于两个流形 $X$ 与 $Y$，镜像对称是导出 Hodge 数等式
%
\begin{equation}
  h^{p,q}(X)=h^{d-p,q}(Y),\qquad d:=\dim_{\C}X=\dim_{\C}Y
  \tag{1}
\end{equation}
%
的对称性。$X$ 与 $Y$ 取连通的代数复流形来考虑，则 $h^{0,0}(X)=1=h^{0,0}(Y)$，因此要使式 (1) 成立，就必须 $h^{dx,0}(X)=1=h^{dy,0}(Y)$ 成立，也就是说 $X$ 与 $Y$ 必须是 Calabi--Yau 流形。

由于同时考虑不同流形 $X$ 与 $Y$ 的组合这件事本身很困难，镜像对称虽可解释为复共轭、Serre 对偶这些自然对称性的延伸，却长时间未被发现。

满足式 (1) 的两个 Calabi--Yau 流形的组合称为拓扑镜像对。那么，拓扑镜像对该怎样构造呢？

% ===== PDF p.147 (原书 p.139) =====
答案之一，是根据奇点论给出的。

\section{Jacobi 环的 Poincaré 多项式}

与第 5 章一样，这里也考虑超曲面奇点，即一个多项式的零点集合中的奇点。所谓由多项式 $f(x_1,\dots,x_n)=0$ 定义的图形的奇点，指的是满足 $f=\partial f/\partial x_1=\cdots=\partial f/\partial x_n=0$ 的点（按重数计入）。特别地，我们考虑这样的情形：Jacobi 环
%
\begin{equation}
  J(f):=\C[x_1,\dots,x_n]\Big/\Big(\frac{\partial f}{\partial x_1},\dots,\frac{\partial f}{\partial x_n}\Big)
  \tag{2}
\end{equation}
%
在 $\C$ 上为有限维；并且，对于某组正整数 $a_1,\dots,a_n$ 与 $h$，由
%
\[
  E:=\sum_{i=1}^{n}\frac{a_i}{h}\cdot x_i\frac{\partial}{\partial x_i}
\]
%
定义的 Euler 向量场 $E$ 满足 $Ef=f$。满足这些条件的奇点称为（带权的）齐次孤立超曲面奇点。例如，对 $A_l$ 型奇点 $f(x):=x^{l+1}$，有 $J(f)=\C[x]/(x^l)$，它在 $\C$ 上是 $l$ 维的，当然有限维；再取 $a:=1$、$h:=l+1$，则% [Tutor Note: 原书此式印作 $\frac{1}{l+1}\cdot\frac{\partial f}{\partial x}=x^{l+1}=f(x)$。按定义 $Ef=\frac{a}{h}x\frac{\partial f}{\partial x}$，应为 $\frac{1}{l+1}\cdot x\cdot\frac{\partial f}{\partial x}=x^{l+1}=f(x)$，原书漏印了一个 $x$。照录原书。]
%
\[
  \frac{1}{l+1}\cdot\frac{\partial f}{\partial x}=x^{l+1}=f(x)
\]
%
所以上述两个条件都被满足。

正整数 $a_1,\dots,a_n$ 是各变量 $x_i$ 的权，$h$ 是多项式 $f(x_1,\dots,x_n)$ 的权，并已归一化，使得用 $E$ 算出的 $f$ 的权为 1。由这些权，Jacobi 环 $J(f)$ 成为分次环。也就是说，它有分解
%
\[
  J(f)=\bigoplus_{d\in\frac{1}{h}\Z_{\ge 0}}J(f)_d,\qquad
  J(f)_d=\{g\in J(f)\mid Eg=dg\}.
\]
%
此时考虑级数
%
\[
  \chi(f,t):=\sum_{d\in\frac{1}{h}\Z_{\ge 0}}\dim_{\C}J(f)_d\cdot t^{d},
\]
%
并称之为 Jacobi 环的 Poincaré 级数。由于假定 Jacobi 环 $J(f)$ 有限维，$\chi(f,t)$ 是 $t^{1/h}$ 的多项式，称为 Poincaré 多项式。特别地，由式 (2)，
% ===== PDF p.148 (原书 p.140) =====
\[
  \chi(f,t)=\prod_{i=1}^{n}\frac{t^{1-\frac{a_i}{h}}-1}{t^{\frac{a_i}{h}}-1}
\]
%
成立，这一点可以简单地证明。而且，由这个式子可知，Poincaré 多项式 $\chi(f,t)$ 具有对偶性
%
\[
  \chi(f,t)=t^{d_f}\cdot\chi(f,t^{-1}),\qquad
  d_f:=\sum_{i=1}^{n}\Big(1-2\frac{a_i}{h}\Big)
\]
%
这一重要性质。用 $A_l$ 型奇点 $f(x)=x^{l+1}$ 的例子来看，
%
\[
  \chi(x^{l+1},t)=1+t^{\frac{1}{l+1}}+\cdots+t^{\frac{l-2}{l+1}}+t^{\frac{l-1}{l+1}},\qquad
  d_{x^{l+1}}=\frac{l-1}{l+1},
\]
%
于是可以确认：在 $\frac{i}{l+1}\leftrightarrow\frac{l-1-i}{l+1}$ 的互换下，Poincaré 多项式不变。

令 $\mu_f:=\chi(f,1)$，把 Poincaré 多项式展开为
%
\[
  \chi(f,t)=\sum_{i=1}^{\mu_f}t^{p_i},\qquad 0=:p_1\le p_2\le\cdots\le p_{\mu}:=d_f,
\]
%
则对偶性可表为 $p_i+p_{\mu_f-i+1}=d_f$。也就是说，可知指数 $p_i$（$i=1,\dots,\mu_f$）的平均值为 $d_f/2$。

指数的平均值既然知道了，再来求指数的方差吧。做一道难度适中的习题式简单计算，就能得到
%
\begin{equation}
  \sum_{i=1}^{\mu_f}\Big(p_i-\frac{d_f}{2}\Big)^2=\frac{d_f}{12}\cdot\mu_f
  \tag{3}
\end{equation}
%
。

\section{不可思议的相似性}

再回到 Calabi--Yau 流形的例子吧。对 Calabi--Yau 流形 $X$，把 Hodge 数的 Poincaré 多项式定义为母函数
%
\[
  \chi(X)(y,\bar y):=\sum_{p,q=0}^{dx}(-1)^{p+q}h^{p,q}(X)y^{p}\bar y^{q}.
\]
%
那么，由 Serre 对偶性，
% ===== PDF p.149 (原书 p.141) =====
\[
  \chi(X)(y,\bar y)=(y\bar y)^{dx}\chi(X)(y^{-1},\bar y^{-1})
\]
%
成立。此外，关于 Hodge 指数的方差，令 $\chi(X):=\chi(X)(1,1)$、$c_k(X)$ 为第 $k$ 个 Chern 类，则已知公式
%
\begin{equation}
  \sum_{p,q=0}^{dx}(-1)^{p+q}h^{p,q}(X)\Big(p-\frac{n}{2}\Big)^2
  =\frac{dx}{12}\chi(X)+\frac{1}{6}\int_X c_1(X)\wedge c_{dx-1}(X)
  \tag{4}
\end{equation}
%
（Libgober--Wood, Borisov）。% [Tutor Note: 式（4）中印作 $(p-n/2)^2$，同页其余处一律用 $dx=\dim_{\C}X$；二者指同一量，属原书记号不统一。照录原文。]


是不是觉得这个式子有点眼熟？Calabi--Yau 流形满足 $c_1(X)=0$。对奇点的方差公式 (3) 与 Calabi--Yau 流形的 Hodge 指数方差公式 (4)，在对应 $p_i\leftrightarrow p$、$d_f\leftrightarrow dx$、$\mu_f\leftrightarrow\chi(X)$ 之下，极其相似呢\footnote{这些指数方差公式，与第 1 章中的 Virasoro 猜想相关，非常重要。此外，对非齐次奇点也能定义指数，并与关于其方差的一般猜想相关。这两者都是非常专门的内容，这里无法涉及。}。

这样一来就明白：在 Hodge 数的层面上，齐次超曲面奇点表现出与 Calabi--Yau 流形非常相似的行为。成问题的是：正如在 $A_l$ 型奇点的例子中所见，应当对应于 Calabi--Yau 流形维度的量 $d_f$、以及应当解释为 Hodge 指数 $p$ 的奇点指数 $p_i$，一般不是整数而是分数。

设 $f_5(x_1,\dots,x_5):=x_1^5+\cdots+x_5^5$，则 $f_5=0$ 在 $\CP^4$ 中确定一个 Calabi--Yau 流形 $X_5$。另一方面，$f_5=0$ 也可以看作 $\C^5$ 中的齐次超曲面奇点。由上面看到的相似性，例如想从多项式 $f_5$ 求出 Calabi--Yau 流形 $X_5$ 的 Hodge 数等等。这会让人想再精确地考察奇点与 Calabi--Yau 流形的关系。实际上可以更精确地比较，而从拓扑镜像对称的观点出发，就会到达下节轨形 Poincaré 多项式的概念。

\section{轨形 Poincaré 多项式}

在前面的例子里，把 $X_5$ 与“把 $f_5=0$ 看作 $\C^5$ 中的奇点”这两者相比，差别在哪里？差别在于：$X_5$ 还是
% ===== PDF p.150 (原书 p.142) =====
$\CP^4$ 还是 $\C^5$，此外还要同时处理作用在多项式与空间上的乘法群 $\Z/5\Z$（或 $\C^*$）——这就是考虑 $X_5$ 的含义。也就是说，$X_5$ 可以看作奇点与群的组 $(f_5,\Z/5\Z)$。这一点我们更具体地来看。

考虑 $GL(n,\C)$ 的有限阿贝尔子群 $G$。特别地，假定 $G$ 的任意元素在 $GL(n,\C)$ 对 $\C^n$ 的自然作用下可写成 $\mathrm{diag}(e[\omega_1\alpha_1],\dots,e[\omega_n\alpha_n])$ 的形式，并具有使 $f$ 不变这一性质。这里 $e[\,\cdot\,]:=\exp(2\pi\sqrt{-1}\,\cdot\,)$，% [Tutor Note: 原书此处印作 $\omega_i:=\alpha_i/h$ 以及 $\alpha_i\in\Z$；但 $\alpha_i$ 同时又是群元素的分量（Vafa 公式中 $\alpha,\beta\in G$ 求和），按第 5 章式（2）的权 $a_i$，此处疑应为 $\omega_i:=a_i/h$。照录原书。]
$\omega_i:=\alpha_i/h$ 以及 $\alpha_i\in\Z$。

在这些准备下，对组 $(f,G)$ 可以给出称为\textbf{轨形 Poincaré 多项式}的不变量 $\chi(f,G)(y,\bar y)$（Vafa 的公式）：
%
\[
\begin{aligned}
  \chi(f,G)(y,\bar y):={}&\frac{(-1)^n}{|G|}\sum_{\alpha,\beta\in G}
  \prod_{\omega_i\alpha_i\notin\Z}(y\bar y)^{\frac{1-2\omega_i}{2}}
  \Big(\frac{y}{\bar y}\Big)^{-\omega_i\alpha_i+[\omega_i\alpha_i]+\frac{1}{2}}\\
  &\times\prod_{\omega_i\alpha_i\in\Z}e\Big[\omega_i\beta_i+\frac{1}{2}\Big]
  \frac{1-e[(1-\omega_i)\beta_i](y\bar y)^{1-\omega_i}}
       {1-e[\omega_i\beta_i](y\bar y)^{\omega_i}}.
\end{aligned}
\]
%
这里，$[\omega_i\alpha_i]$ 表示不超过 $\omega_i\alpha_i$ 的最大整数。令 $t=y\bar y$，则对 $f$ 与平凡群 $G=\{1\}$ 的组，轨形 Poincaré 多项式为
%
\[
  \chi(f,\{1\})(y,\bar y)=\chi(f,t),
\]
%
这与先前定义的 Poincaré 多项式一致。

从带群作用的超曲面奇点、即组 $(f,G)$ 得到种种不变量，这在物理中称为\textbf{Landau--Ginzburg 轨形模型}（Landau--Ginzburg orbifold model）并得到研究。从历史上看，正是从这种 Landau--Ginzburg 轨形模型的拓扑镜像对称研究出发，发展出了 Calabi--Yau 流形的镜像对称研究。

\section{奇点的拓扑镜像对称}

把式 (1) 用 Poincaré 多项式改写，则
%
\[
  \chi(X)(y,\bar y)=(-1)^{d_Y}y^{d_Y}\chi(Y)(y^{-1},\bar y)
\]
%
成立。因此，可以这样定义：组 $(f,G)$ 与组 $(f^*,G^*)$ 互为拓扑镜像对，是指
%
\[
  \chi(f,G)(y,\bar y)=(-1)^{d_{f^*}}y^{d_{f^*}}\chi(f^*,G^*)(y^{-1},\bar y)
\]
%
成立。哪些组 $(f,G)$ 与组 $(f^*,G^*)$
% ===== PDF p.151 (原书 p.143) =====
在什么条件下互为拓扑镜像对，部分物理想法虽已知道，但要以一般形式在数学上证明，是非常困难的问题。不过，组成拓扑镜像对的组 $(f,G)$ 与组 $(f^*,G^*)$ 的例子已被发现非常多，有很多应用。

再次考虑多项式 $f_5(x_1,\dots,x_5)=x_1^5+\cdots+x_5^5$。令 $G$ 为由 $\mathrm{diag}(e[\omega_1],\dots,e[\omega_n])$ 生成的 $\Z/5\Z$，$f_5^*:=f_5$，$G^*$ 为满足前节条件、且使单项式 $x_1\cdots x_5$ 不变的最大有限阿贝尔群。设 $X_5$ 为由 $\{f_5=0\}\subset\CP^4$ 给出的 Calabi--Yau 流形，$X_5^*$ 为把 $X_5$ 按群 $G^*/G$ 取商所得空间的某个极小消解，则
%
\[
  \chi(f,G)(y,\bar y)=\chi(X_5)(y,\bar y),\qquad
  \chi(f^*,G^*)(y,\bar y)=\chi(X_5^*)(y,\bar y),
\]
%
特别地，组 $(f_5,G)$ 与组 $(f_5^*,G^*)$ 成为拓扑镜像对。这就是著名的 5 次超曲面的拓扑镜像流形的构成法。用同样的方法，构成了 7555 个互为拓扑镜像对的 3 维 Calabi--Yau 超曲面。

\section{Milnor 纤维的几何学}

此前我们关注的是拓扑镜像对称。为讲述与同调镜像对称相关的话题，下面就复 2 维情形简单触及从超曲面奇点得到的几何不变量。设 $f(x_1,x_2,x_3)$ 为 2 维孤立超曲面奇点，则
%
\begin{equation}
  f:\C^3\setminus f^{-1}(0)\longrightarrow\C\setminus\{0\}
  \tag{5}
\end{equation}
%
是拓扑局部平凡的纤维束。于是，考虑称为 \textbf{Milnor 纤维}的一般纤维 $X_1:=f^{-1}(1)$。$X_1$ 是没有奇点的复流形。$X_1$ 是复 2 维（实 4 维）的，因此由 Poincaré 对偶，存在称为\textbf{相交形式}的对称双线性形式
%
\begin{equation}
  I:H_2(X_1,\Z)\times H_2(X_1,\Z)\longrightarrow\Z.
  \tag{6}
\end{equation}
%
由 Milnor 定理可知，$H_2(X_1,\Z)$ 是秩 $\mu_f$ 的自由加群，作为其生成元，可以取塌缩到奇点的、称为\textbf{消没闭链}的子流形。

% ===== PDF p.152 (原书 p.144) =====
其中，组 $(H_2(X_1,\Z),-I)$ 称为 Milnor 格（lattice）。第 5 章中令 $f=x^{l+1}$ 时，看到 0 维球面\footnote{本例中，设 $\omega$ 为 $1$ 的本原 $l+1$ 次单位根时，表示由形如 $[\omega^i]-[\omega^j]$ 的两点所给出的图形。}塌缩到原点的情形，还记得吗？这里所述的内容是把它稍稍一般化。% [Tutor Note: 原书此处作「原始 $l+1$ 次単位根」，即本原（primitive）单位根；「原始」在本译本中固定对应 primitive form（原始形式），故正文改用「本原」以免混淆。]

又，由式 (5) 可导出群的表示
%
\[
  \rho:\pi_1(\C\setminus\{0\},*)\;(\cong\Z)\longrightarrow
  \mathrm{Aut}_{\Z}(H_2(X_1,\Z),I).
\]
%
$\Z$ 的生成元 $1$ 的表示 $c:=\rho(1)$ 称为\textbf{Milnor 单值性}\footnote{在 $f=x^{l+1}$ 的例子中，Milnor 单值性 $c$ 是把 0 维球面 $[\omega^i]-[\omega^j]$ 映为 $[\omega^{i+1}]-[\omega^{j+1}]$ 的变换。}。把消没闭链构成的集合记作 $\mathcal{R}$，则四元组 $(H_2(X_1,\Z),-I,\mathcal{R},c)$ 满足“\textbf{一般根系}”这一古典根系扩张概念的公理（斋藤恭司）。特别地，若从 ADE 型（含 $A_l$ 型奇点）的（带权）齐次孤立超曲面奇点出发，作为一般根系，可得到与该型相应的古典根系。

\section{奇点的同调镜像对称}

在 $\C^3$ 上赋予标准辛形式时，关于它诱导的 $X_1$ 上的辛形式，消没闭链可以取成特殊的子流形，即\textbf{Lagrangian 子流形}。于是，以 Milnor 纤维 $X_1$ 的 Lagrangian 子流形为对象、以它们的 Floer 同调为态射空间，伴随 $f$ 的 Fukaya 范畴 $D^bLag(f)$ 便有望自然地构成。不过首先有几个问题：Milnor 纤维 $X_1$ 非紧；对象类最好从 Lagrangian 子流形再稍加扩大，等等。所以“伴随 $f$ 的 Fukaya 范畴 $D^bLag(f)$”这一写法本身，从数学上正确定义它开始就是难题。一旦这些问题解决，前节给出的一般根系 $(H_2(X_1,\Z),-I,\mathcal{R},c)$ 就会作为 Fukaya 范畴 $D^bLag(f)$ 的不变量导出。

所谓\textbf{同调镜像对称}，是指把由辛几何确定的 $A_\infty$ 范畴（或三角范畴）与由复几何确定的相应范畴互换的对称性。人们期待它是超越 Calabi--Yau 流形、在非常大的类上成立的对称性。特别地，拓扑镜像对称与经典镜像对称，被认为应是由更强的同调镜像对称导出的推论。

% ===== PDF p.153 (原书 p.145) =====
带群作用的孤立奇点之间存在拓扑镜像对称
%
\[
  (f,\{1\})\Longleftrightarrow(f^*,G^*)
\]
%
\footnote{这里 $G^*$ 取为使 $f^*$ 不变（满足第 5 节条件）的 $GL(3,\C)$ 中的最大有限阿贝尔群。}。因此，如果拓扑镜像对称被范畴化为同调镜像对称，那么对于奇点 $f$，组 $(f^*,G^*)$ 便确定。通过 $A_\infty$ 范畴可以构成由代数方式确定的三角范畴 $\mathcal{T}_{f^*,G^*}$，并期待它与 Fukaya 范畴 $D^bLag(f)$ 作为三角范畴等价。\textbf{这就是奇点的同调镜像对称}\footnote{奇点的同调镜像对称在一般维数上也可期待。为简化说明，以下仍限定在 2 维。}。更一般地，把 $(f,G)$ 与 $(f^*,G^*)$ 取为拓扑镜像对时，期待三角范畴的等价
%
\[
  D^bLag(f,G)\cong\mathcal{T}_{f^*,G^*},
\]
%
这便是一般情形下的奇点同调镜像对称猜想\footnote{当然，与通常的同调镜像对称猜想一样，为了让“等价”一词有意义，从定义两边的范畴开始就已经是问题（猜想）的一部分。不过，期待它的数学定理、背景、现象、根据等，已在 1980 年代的奇点论中被发现。}。

用一般根系 $(H_2(X_1,\Z),-I,\mathcal{R},c)$ 更具体地说，下面的对应关系可期待由同调镜像对称成立。

下页表格中的两侧，可以解释为同一数学对象的不同实现。作为镜像对称的应用，例如有理曲线计数等，把辛几何中非常困难的不变量用复几何中可计算的量来研究，这样的做法一直在进行。奇点的情形也同样，通过研究由代数方式确定的三角范畴 $\mathcal{T}_{f^*,G^*}$，就有可能研究极为超越的对象——消没闭链的几何。

% ===== PDF p.154 (原书 p.146) =====
\begin{table}[htbp]
\centering
  \tabtitle{1}{对应关系（原书 p.146）}
\begin{tabular}{c|c}
\hline
辛几何 & 复几何 \\
\hline
$(f,\{1\})$ & $(f^*,G^*)$ \\
\hline
$H_2(X_1,\Z)$\ 中间同调群 & $K_0(\mathcal{T}_{f^*,G^*})$\ Grothendieck 群 \\
\hline
$-I$，$I$：相交形式 & $\chi+{}^t\chi$，$\chi$：Euler 形式 \\
\hline
$\mathcal{R}$\ 消没闭链的类 & 不可分解对象的类 $\in K_0(\mathcal{T}_{f^*,G^*})$ \\
\hline
$c$\ Milnor 单值性 & $-[S]\in\mathrm{Aut}(K_0(\mathcal{T}_{f^*,G^*}))$（$S$：Serre 函子）\\
\hline
\end{tabular}
\end{table}

\section{非交换几何的登场与根系的实现}

$\mathcal{T}_{f^*,G^*}$ 是如何定义的，下面连同背景稍作说明。组 $\mathcal{A}_{f^*}:=(\C[x_1,x_2,x_3],f^*)$ 可以看作仅由 1 个对象组成的 $A_\infty$ 范畴（$A_\infty$ 代数）。用口号式的话说，就是把可换奇点重新看作以可换环 $\mathcal{A}_{f^*}$ 为函数环的非交换空间。然后，考虑 $\mathcal{A}_{f^*}$ 上的 $\widehat{G}^*$-等变右 $A_\infty$-加群构成的三角范畴
%
\[
  \mathcal{T}_{f^*,G^*}:=D_{\widehat{G}^*}(\mathrm{Mod}\text{-}\mathcal{A}_{f^*}).
\]
%
这里 $\widehat{G}^*$ 是群 $G^*$ 由 $\Z$ 的扩张，其来源是多项式 $f^*$ 带有 $\C^*$-作用。

这个范畴，是代数流形 $X$ 上群 $G$ 作用时，$G$-等变的拟凝聚层（quasi-coherent sheaf）的复形构成的导出范畴 $D_G(\mathrm{Qcoh}\,X)$ 的类似物。$X$ 光滑，且 $X/G$ 具有保持典范丛的奇点消解\footnote{因此，是指满足第一 Chern 类 $c_1(Y)=c_1(X/G)=0$ 的奇点消解。} $Y$ 时，$D_G(\mathrm{Qcoh}\,X)$ 的紧对象\footnote{它对应的余可表函子与无限直和可交换的对象。}全体构成的满三角子范畴，与 $G$-等变凝聚层的有界复形构成的导出范畴 $D^b_G(\mathrm{coh}\,X)$ 等价，特别地
%
\[
  D^b_G(\mathrm{coh}\,X)\cong D^b(\mathrm{coh}\,Y)
\]
%
成立。如第 10 章所述，凝聚层的有界导出范畴 $D^b(\mathrm{coh}\,Y)$ 可看作复几何侧的 D-膜范畴，在镜像对称中应与辛几何侧的 Fukaya 范畴等价。

% ===== PDF p.155 (原书 p.147) =====
下面举 1 个具体例子。对于第 6 节登场的 $\CP^4$ 的 5 次超曲面 $X_5$ 与其拓扑镜像流形 $X_5^*$，
%
\[
  D^b_{G^*/G_0}(\mathrm{coh}\,X_5)\cong D^b(\mathrm{coh}\,X_5^*)
\]
%
成立。这可以看作第 6 节拓扑镜像对构造的范畴化；特别地，Poincaré 多项式的等式可由这个三角范畴的等价导出。

现在考虑的奇点、即非交换空间 $\mathcal{A}_{f^*}$，可以说具有“非交换 Calabi--Yau 流形”这样非常良好的性质。并且可以认为“商空间 $\mathcal{A}_{f^*}/G^*$ 具有保持典范丛的非交换奇点消解”。基于可换情形的类比，$\mathcal{T}_{f^*,G^*}$ 的紧对象全体构成的满三角子范畴就是所求的三角范畴 $\mathcal{T}_{f^*,G^*}$。而 $\mathcal{T}_{f^*,G^*}$ 与 $A_\infty$ 范畴 $\mathcal{A}_{f^*}$ 的 $\widehat{G}^*$-等变导出范畴 $D_{\widehat{G}^*}(\mathcal{A}_{f^*})$ 等价，可以这样认为。令人惊讶的是，这个三角范畴 $D_{\widehat{G}^*}(\mathcal{A}_{f^*})$ 与某个三角范畴一致。后者由带群作用的矩阵因子化构成，是 1980 年代中期在奇点及极大 Cohen--Macaulay 模的研究中引入的。

在辛几何侧、消没闭链的几何侧，已知可以取到消没闭链的一组 distinguished basis（特别基），它成为根格 $(H_2(X_1,\Z),-I)$ 的“好基”。它的范畴论对应物，是称为 exceptional collection 的三角范畴的对象序列 $(\mathcal{E}_1,\dots,\mathcal{E}_\mu)$。对于好的一类齐次奇点 $f$，可以证明 $D_{\widehat{G}^*}(\mathcal{A}_{f^*})$ 由 exceptional collection 生成。特别地，可知 Grothendieck 群 $K_0(\mathcal{T}_{f^*,G^*})$ 的好基可以取到。

进一步，若 $\mathcal{E}_i$ 之间除 0 次以外的所有 Ext 群都为零，则自同态代数 $E:=\mathrm{End}(\mathcal{E}_1\oplus\cdots\oplus\mathcal{E}_\mu)$ 是非交换有限维代数。特别地，它等于某个箭图的路径代数按关系式取商所得。而有限生成右 $E$-模的导出范畴 $D^b(\mathrm{mod}\text{-}E)$ 与 $D_{\widehat{G}^*}(\mathcal{A}_{f^*})$ 作为三角范畴等价，因此奇点 $f$ 的镜像流形，就是函数环为 $E$ 的“非交换流形”。又，正如拓扑镜像对称所期待，它也成为“具有分数维数 $d_f$ 的非交换 Calabi--Yau 流形”。

% ===== PDF p.156 (原书 p.148) =====
例如，若 $f$ 为 ADE 型奇点，则 $E$ 是对应型的 Dynkin 图适当地赋予定向后所得 Dynkin 箭图的路径代数。由此，从三角范畴 $\mathcal{T}_{f^*,G^*}$ 经 Grothendieck 群，可以得到古典根系；这为奇点也存在同调镜像对称提供了根据。

\section{结语}

第 5 章与本章的解说想传达的氛围，是这些乍看不同的对象之间所存在的不可思议的关联：Calabi--Yau 流形、奇点、根系、平坦结构、量子上同调。此外还有不少有趣的话题：奇点的范畴（Buchweitz、Orlov）与矩阵因子化的范畴之间的等价，Bridgeland--Douglas 的 $\Pi$-稳定性与周期映射，通过 $\mathcal{T}_{f,\{1\}}$ 的“$A_\infty$-变形”导出超曲面奇点的平坦结构，等等。遗憾的是，许多有趣的话题只能省略，希望有机会再作解说。

% 恢复全书默认的罗马数字节号（本章原书为阿拉伯数字，见开头说明）。
\renewcommand{\thesection}{\Roman{section}}
